\quad Os algoritmos de ordenação apresentam grande importância no estudo teórico de algoritmos, bem como em aplicações práticas do cotidiano. Esses algoritmos são responsáveis por organizar conjuntos de elementos de acordo com determinado critério, facilitando a busca, a visualização e o processamento dos dados. Diante da existência de diferentes métodos de ordenação, torna-se importante compreender seu funcionamento e analisar seus custos para identificar suas características e limitações.

\quad Neste trabalho, são analisados seis algoritmos de ordenação: \textit{Insertion Sort}, \textit{Bubble Sort}, \textit{Selection Sort}, \textit{Shell Sort}, \textit{Merge Sort} e \textit{Quick Sort}. Cada algoritmo apresenta uma estratégia diferente para organizar os elementos de um vetor. Enquanto o \textit{Insertion Sort} realiza inserções progressivas, o \textit{Bubble Sort} utiliza trocas entre elementos adjacentes, o \textit{Selection Sort} seleciona sucessivamente o menor elemento, o \textit{Shell Sort} utiliza intervalos (\textit{gaps}) para realizar ordenações parciais, o \textit{Merge Sort} emprega a estratégia de divisão e conquista com intercalação de subvetores e o \textit{Quick Sort} utiliza divisão e conquista com particionamento em torno de um pivô. Para o \textit{Quick Sort}, foram implementadas duas versões: uma que utiliza o primeiro elemento como pivô e outra que define o pivô por meio da média aritmética dos elementos da partição.

\quad Para cada algoritmo, são apresentados sua implementação, um teste de mesa manual e uma análise de complexidade, considerando os diferentes casos de execução. Também são realizados experimentos com entradas ordenadas de forma crescente, decrescente e aleatória, utilizando diferentes tamanhos de instância. Os tempos de execução obtidos são apresentados por meio de tabelas e gráficos, permitindo relacionar o comportamento observado experimentalmente com a análise teórica de cada algoritmo, inclusive a influência da escolha do pivô no desempenho do \textit{Quick Sort}.

\quad Dessa forma, o trabalho busca compreender as diferenças entre os algoritmos estudados e analisar como suas características influenciam o desempenho durante a ordenação de diferentes tipos e tamanhos de entrada.